% Formato del título de capítulos y secciones
% \titleformat{\chapter}[block]{\titlerule[0.8pt]\normalfont\huge\bfseries}{\thechapter.}{.5em}{\Huge}[{\titlerule[0.8pt]}]
% \titlespacing*{\chapter}{0pt}{-19pt}{25pt}
% \titleformat{\section}[block]{\normalfont\Large\bfseries}{\thesection.}{.5em}{\Large}

\titleformat{\chapter}[block]{\normalfont\huge\bfseries}{\thechapter.}{.5em}{\Huge}[{\titlerule[0.8pt]}]

\titlespacing*{\chapter}{0pt}{-19pt}{25pt}
\titleformat{\section}[block]{\normalfont\Large\bfseries}{\thesection.}{.5em}{\Large}



% Formato del código fuente con lstlisting
\lstset{
  basicstyle=\ttfamily,
  breaklines=true,
}

% Márgenes
\geometry{
    a4paper,
    margin=2.5cm
}

% Limite de profundidad del índice
\setcounter{tocdepth}{2}

% Indentación de párrafos
\setlength{\parindent}{0.5cm}

% Espaciado entre párrfos
\setlength{\parskip}{2mm}


% Espaciado entre líneas (equivale a 1.5 en Word según https://tex.stackexchange.com/questions/65849/confusion-onehalfspacing-vs-spacing-vs-word-vs-the-world
\linespread{1.25}
% \linespread{1.0} lo dejo así hasta que sea la versión final

\renewcommand{\lstlistingname}{Extracto de código}
\renewcommand*{\lstlistlistingname}{Índice de extractos de código}

% Label de las referencias
\crefname{figure}{Fig.}{Figs.} % Singular y plural para figuras
\crefname{table}{Table}{Tables} % Singular y plural para tablas
\crefname{equation}{Eq.}{Eqs.} % Singular y plural para ecuaciones
\crefname{chapter}{Capítulo}{Capítulos} % Singular y plural para 


%------------ Caption debajo de las tablas
\DeclareCaptionLabelSeparator*{spaced}{\\[1.5ex]}

% \captionsetup[table]{textfont=normalsize,format=plain,justification=justified,
%   singlelinecheck=false,labelsep=period,skip=0.5pt}

%--------------- Caption de tablas y figuras (edit Iván)
\captionsetup[table]{font={small, singlespacing},
  format=plain, justification=justified,
  singlelinecheck=false, labelsep=period, skip=2mm}

\captionsetup[figure]{font={small, singlespacing},
  format=plain, justification=justified,
  singlelinecheck=false, labelsep=period, skip=2mm}

% Tipo de columnas con alineación y ancho personalizado
% (L = left, C = centred)
\newcolumntype{C}[1]{>{\centering\arraybackslash}p{#1}}
\newcolumntype{L}[1]{>{\raggedright\arraybackslash}p{#1}}

% Configuración del idioma
\selectlanguage{spanish}
\decimalpoint % Fuerza que los números en math mode usen el punto decimal

% Que el espacio en blanco quede al final de la página, no uniforme
\raggedbottom


% Evitar saltos de página justo abajo de títulos de section y subsection
\makeatletter
\renewcommand{\section}{\@startsection{section}{1}{0pt}{10pt}{10pt}{\normalfont\Large\bfseries\@afterheading}}
\renewcommand{\subsection}{\@startsection{subsection}{2}{0pt}{10pt}{10pt}{\normalfont\large\bfseries\@afterheading}}
\renewcommand{\subsubsection}{\@startsection{subsubsection}{3}{0pt}{8pt}{8pt}{\normalfont\normalsize\bfseries\@afterheading}}
\makeatother

% comando para cpp
\newcommand\cpp{C\nolinebreak[4]\hspace{-.05em}\raisebox{.4ex}{\relsize{-3}{\textbf{++}}}\xspace}


% Configurar hipervínculos en el texto
\hypersetup{
    colorlinks=true, % Hace que los enlaces sean de color en vez de cajas
    linkcolor=black, % Color de enlaces internos (índice, referencias)
    citecolor=blue, % Color de referencias bibliográficas
    filecolor=blue, % Color de enlaces a archivos
    urlcolor=black, % Mantiene las URLs en azul
    pdfborder={0 0 0} % Elimina subrayado en la bibliografía
}


% Ubicar flotantes en el comienzo de la página 
\makeatletter
\setlength{\@fptop}{0pt}
\makeatother

% Que las figuras y tablas sean del mismo tipo, así 
% aparecen en orden
\makeatletter
\let\ftype@table\ftype@figure
\makeatother
